\chapter{Conclusion and Future Work}\label{ch:conclusion}

This thesis presented the design, implementation, and evaluation of NELS, a ComfyUI plugin for exploring latent spaces in generative diffusion models. The work covered the complete exploration pipeline, including dataset construction, embedding projection, clustering, and a range of navigation operations. Their empirical evaluation examined structural and perceptual behavior, stability, interpretability, and controllability. The results show that latent-space navigation can be made systematic and inspectable, while also revealing that several operations behave less predictably and reliably than their design initially suggests.

These findings point to several directions for future work. The most immediate is a controlled user study examining whether the limitations identified through automated evaluation are also perceived by users and whether the plugin provides an effective environment for latent-space exploration. Further work could refine the coarse variation-strength control of Consistency Check, improve the interpretability of the automatically derived second axis in Semantic Axis Projection through additional constraints or alternative discovery methods, investigate other dimensionality-reduction techniques for more effective visualization, and repeat the evaluation across different diffusion models and larger prompt datasets.

Taken together, NELS provides both a working exploration tool and a first empirical account of how latent-space navigation actually behaves, laying the groundwork for the deeper investigations into latent spaces that future work should pursue.

The complete source code is publicly available in the project repository \cite{davidgh2026nels}.
